\documentclass[11pt,a4paper]{article}

\usepackage{../shared/cathedral-preamble}
\usepackage[table]{xcolor}  % for \rowcolor

\title{%
  Cathedral Next:\\
  Open Horizons in Machine-Verified\\
  Analytic Number Theory%
}
\author{
  Cathedral Project Notes
}
\date{April 2026}

\begin{document}

\maketitle

\begin{abstract}
We survey the current state of the Cathedral---a Lean~4 formalization
reducing the Riemann Hypothesis to six machine-checked axioms via the
Nyman--Beurling criterion---and chart directions for future work. We
identify three horizons: (1) the \emph{axiom horizon}, where the
remaining six crown axioms yield to formalization (one has already
fallen since the initial survey); (2) the
\emph{generalization horizon}, where the framework extends to Dirichlet
$L$-functions and the Generalized Riemann Hypothesis; and (3) the
\emph{interpretation horizon}, where the Liouville parity symmetry,
octonionic partition, and spectral gap phenomena discovered during
formalization connect to mathematical physics and random matrix theory.

This document is not a claim of proof. It is a map of what the Cathedral
revealed, and where the trails lead.

\textbf{Note added April~22, 2026.} Since the initial version of this
paper (April~20), Campaign~B (the Abel Engine) has been completed:
\lean{abel\_mertens\_tail\_raw} is now a \textbf{THEOREM}. The Perron
formula infrastructure has been fully proved. The Borel--Carath\'eodory
theorem has been discovered in Mathlib, opening an unexpected route to
Campaign~C. The axiom count on the crown path has decreased from 7 to~6.
\end{abstract}

\tableofcontents

% ================================================================
\section{Where We Stand}
\label{sec:status}
% ================================================================

\subsection{The Crown Theorem}

The Cathedral's crown theorem \lean{nyman\_beurling\_equivalence}
establishes:
\[
  \RH \;\iff\;
  \forall\,\varepsilon > 0,\;\exists\,N_0,\;\forall\,N \geq N_0,\;
  \exists\,v : \mathrm{Fin}(N-1) \to \RR,\quad
  \int_0^1 \bigl(1 - f_{N,v}(x)\bigr)^2\,dx < \varepsilon
\]
where $f_{N,v}(x) = \sum_{k=2}^{N} v_k\,\{1/(kx)\}$ is a linear
combination of B\'aez-Duarte basis functions.

The proof decomposes into two pillars:
\begin{itemize}
  \item \textbf{Pillar~I (Converse)}: $d_N^2 \to 0 \implies \RH$.
    Proved via the Rank-1 Mellin Miracle. \emph{Zero custom axioms.}
  \item \textbf{Pillar~II (Forward)}: $\RH \implies d_N^2 \to 0$.
    Proved via the Final Dragon chain. \emph{Six custom axioms.}
    (Reduced from seven on April~22, 2026, when the Abel tail axiom
    was graduated to a theorem.)
\end{itemize}

\subsection{The Six Crown Axioms}

\begin{table}[h]
\centering
\small
\begin{tabular}{@{}clll@{}}
\toprule
\textbf{\#} & \textbf{Axiom} & \textbf{Content} & \textbf{Tier} \\
\midrule
1 & \lean{rh\_implies\_mertens\_bound} & $\RH \implies |M(x)| = O(x^{1/2}\log^2 x)$ & 1 (RH) \\
2 & \lean{pnt\_mu\_div\_k} & $\sum \mu(k)/k \to 0$ & 2 (PNT) \\
3 & \lean{pnt\_mu\_log\_div\_k} & $\sum \mu(k)\ln(k)/k \to -1$ & 2 (PNT) \\
4 & \lean{pnt\_mu\_log\_sq\_div\_k} & $\sum \mu(k)\ln^2(k)/k \to -2\gamma$ & 2 (PNT) \\
\rowcolor{green!10}
$\times$ & \lean{abel\_mertens\_tail\_raw} & Abel summation tail bounds & \textbf{GRADUATED} \\
5 & \lean{millennium\_covariance\_cancellation} & 2D covariance bound & 3 (Parseval) \\
6 & \lean{vasyunin\_offdiag\_integral} & Off-diagonal Gram identity & 3 (Vasyunin) \\
\bottomrule
\end{tabular}
\caption{The crown axioms. Originally seven; reduced to six on April~22, 2026
when the Abel tail axiom was graduated to a theorem. Tier~1 encodes RH
content; Tier~2 is unconditional (PNT-level); Tier~3 is classical analysis.
The converse direction uses zero custom axioms.}
\label{tab:axioms}
\end{table}

\subsection{What Was Proved Along the Way}

The formalization is not merely a scaffolding of axioms. Among the fully
machine-verified results:

\begin{itemize}
  \item The \emph{Rank-1 Mellin Miracle}: $\mathcal{M}[h_k](\rho) =
    1/(k(\rho-1))$ at zeta zeros (BDMellin.lean, 1066 lines).
  \item The base-case Mellin identity via the Identity Theorem for
    holomorphic functions (IdentityBypass.lean).
  \item The completed zeta bound $\re(\Lambda_0(s)) < 4$ for real
    $s \in (0,1)$ via Jacobi theta kernel estimates (ThetaBound.lean).
  \item No real zeros of $\zeta(s)$ in $(0,1)$ (corollary of the above).
  \item Under RH, $\zeta(s) \neq 0$ for $\re(s) > 1/2$ and $s \neq 1$
    (ZetaConvexity.lean, zero axioms).
  \item $1/\zeta(s)$ is differentiable for $\re(s) > 1/2$ under RH
    (ZetaConvexity.lean, zero axioms).
  \item Weyl's eigenvalue inequality for Hermitian matrix sums
    (RayleighBridge.lean).
  \item The Gram matrix is Hermitian, with spectral decomposition and
    Rayleigh quotient characterization (Defs.lean, ClassRestriction.lean).
  \item The octonionic gap dominance $\lambda_{\min}(G) \le
    \lambda_{\min}(G^{\mathbb{O}})$ via class-restricted Rayleigh
    quotients (ClassRestriction.lean, 647 lines, zero axioms).
  \item The eigenvalue limit $\lim_{N \to \infty} \lambda_{\min}(G_N)$
    exists (MainChain.lean, via monotone convergence).
  \item The Liouville function squares to 1 (Defs.lean, \lean{native\_decide}).
  \item \textbf{[April~22]} Abel tail bounds: $|S_1| \le C \cdot N^{-1/4}$,
    $|S_2+1| \le C \cdot N^{-1/4} \ln N$, $|S_3+2\gamma|$ bounded
    (AbelTail/*.lean, zero axioms --- \emph{formerly Axiom \#5}).
  \item \textbf{[April~22]} Perron kernel bounds and contour vanishing
    (PerronKernel.lean, Perron/Rectangle.lean, zero axioms).
  \item \textbf{[April~22]} The Dirichlet series $1/\zeta(s) =
    \sum \mu(n)/n^s$ for $\re(s) > 1$ (DirichletZetaInverse.lean, zero axioms).
\end{itemize}

% ================================================================
\section{Horizon I: Closing the Axioms}
\label{sec:axioms}
% ================================================================

The six crown axioms partition into three natural campaigns, ordered
by increasing mathematical depth. \textbf{Campaign~B has been completed}
(April~22, 2026).

\subsection{Campaign A: The PNT Axioms (Tier~2)}

\begin{axiombox}[PNT Partial Sums]
\label{ax:pnt}
The three PNT axioms state:
\begin{align}
  \sum_{k=1}^{N} \frac{\mu(k)}{k} &\to 0, \label{eq:pnt1} \\
  \sum_{k=1}^{N} \frac{\mu(k)\ln k}{k} &\to -1, \label{eq:pnt2} \\
  \sum_{k=1}^{N} \frac{\mu(k)\ln^2 k}{k} &\to -2\gamma. \label{eq:pnt3}
\end{align}
\end{axiombox}

These are unconditional consequences of the Prime Number Theorem.
The proof route is well-understood:

\begin{enumerate}
  \item \eqref{eq:pnt1} follows from PNT via
    $1/\zeta(s) = \sum \mu(n)/n^s$ and Abel's limit theorem as $s \to 1^+$.
  \item \eqref{eq:pnt2} follows by differentiating $-\zeta'(s)/\zeta(s)^2$
    and applying the same limit.
  \item \eqref{eq:pnt3} follows from the Laurent expansion of $1/\zeta(s)$
    near $s = 1$.
\end{enumerate}

\begin{direction}[PNT in Lean~4]
Mathlib already contains \lean{riemannZeta} and its analytic continuation.
The obstacle is the Wiener--Ikehara Tauberian theorem, which converts
Dirichlet series asymptotics into partial sum asymptotics. Three approaches:

\begin{enumerate}
  \item \textbf{Formalizing Wiener--Ikehara directly.} This requires
    convolution theory, Fej\'er kernel estimates, and the non-vanishing
    $\zeta(1+it) \neq 0$. The non-vanishing is proved in Mathlib
    (\lean{riemannZeta\_ne\_zero\_of\_one\_le\_re}).

  \item \textbf{The Newman--Korevaar shortcut.} A simplified proof of
    PNT via contour integration, avoiding the Tauberian theorem entirely.
    Requires only the analyticity of $\sum \mu(n)/n^s - 1/(s-1)\zeta(s)$
    in $\re(s) \geq 1$ and a Laplace transform lemma.

  \item \textbf{Selberg--Erd\H{o}s elementary proof.} Entirely real-variable,
    avoiding complex analysis. Longer but requires no Tauberian machinery.
    May be the most Lean-friendly route.
\end{enumerate}

The author's recommendation: the Newman--Korevaar route, since most of the
complex analysis infrastructure already exists in Mathlib. The key lemma
(analytic continuation of $\zeta'(s)/\zeta(s)$ to $\re(s) \geq 1$) is
essentially the non-vanishing on the line, which is already formalized.
\end{direction}

\subsection{Campaign B: The Abel Engine (COMPLETED)}
\label{sec:abel-completed}

\begin{axiombox}[Abel--Mertens Tail --- GRADUATED April~22, 2026]
\lean{abel\_mertens\_tail\_raw}: Given $|M(x)| \le C_m x^{3/4}$ and
the PNT convergences, the Abel summation on the tail
$\sum_{k > N} M(k)/(k(k+1))$ gives:
\[
  |S_1(N)| \le C \cdot N^{-1/4}, \quad
  |S_2(N)+1| \le C \cdot N^{-1/4}\ln N, \quad
  |S_3(N)+2\gamma| \le C \cdot N^{-1/4}\ln^2 N.
\]
\end{axiombox}

\textbf{This axiom is now a THEOREM.} The proof, completed April~22, 2026,
uses a three-term decomposition ($S_1$, $S_2$, $S_3$) with separate
decay arguments for each term:
\begin{enumerate}
  \item $S_1$ decay: Abel summation by parts with integral comparison
    $\sum_{k \geq N} k^{-5/4} \le 4(N-1)^{-1/4}$.
  \item $S_2$ decay: Similar, with an additional $\ln k$ factor absorbed
    by $\ln N$.
  \item $S_3$ bound: Direct from the PNT axiom $S_3 \to -2\gamma$.
\end{enumerate}
The proof required 4 new files in \lean{AbelTail/} (Telescoping, AbelInterior,
LogTailBound, MononotoneBound) plus assembly in \lean{CalcBounds.lean}.
All files have zero sorry.

This graduation reduces the crown path from 7 to 6 axioms and proves
that the ``thermodynamic hierarchy'' of the Mertens function---the
magnetization ($S_1$), susceptibility ($S_2$), and heat capacity ($S_3$)
of the prime number gas---converges at rate $N^{-1/4}$.

\subsection{Campaign C: The Deep Axioms (Tier~3, Axioms~6--7)}

These are the mathematically deepest axioms and will likely be the last
to fall. However, the recent discovery of the Borel--Carath\'eodory
theorem in Mathlib (\texttt{Analysis.Complex.BorelCaratheodory}) opens
unexpected possibilities.

\subsubsection{The Millennium Wall (Axiom~6)}

\begin{axiombox}[Covariance Cancellation]
\lean{millennium\_covariance\_cancellation}: The covariance matrix
$C = G - bb^T$ satisfies $v^T C v \le K_{\mathrm{cov}}/\ln N$ for the
M\"obius witness $v$.
\end{axiombox}

This is the \emph{irreducible analytic content} of the forward direction.
It requires:
\begin{enumerate}
  \item The Mellin--Plancherel isometry converting the bilinear form
    to a critical-line integral.
  \item The Montgomery--Vaughan mean value theorem bounding
    $\int |\zeta(1/2+it) W_N(1/2+it)|^2 / (1/4+t^2)\,dt$.
  \item Contour shifting from $\re(s) = 1/2$ to $\re(s) = 2$,
    picking up the residue at the double pole $s = 1$.
\end{enumerate}

\begin{direction}[The Contour Shift Route]
The Parseval Bridge (cathedral.tex, \S4) shows that the Plancherel
isometry converts the spatial bilinear form into a frequency-domain
integral. The Cathedral already contains the algebraic three-term
decomposition
\[
  \int_0^1 |1 - f_{N,v}|^2\,dx
  = \frac{1}{2\pi}\int_{-\infty}^{\infty}
    \frac{|1 - \zeta(s)W_N(s)|^2}{|s|^2}\,dt,
    \quad s = \tfrac{1}{2}+it.
\]
The contour shift $\re(s) = 1/2 \to \re(s) = 2$ is standard in analytic
number theory (Phragm\'en--Lindel\"of for the horizontal segments,
Cauchy's theorem for the rectangle). The pole at $s = 1$ contributes
a residue involving $W_N(1)$ and $W_N'(1)$, which encode the
$\ln\ln N/\ln N$ correction.

This is probably a 500--1000 line Lean formalization, requiring:
\begin{itemize}
  \item Cauchy's integral formula for meromorphic functions (partially in Mathlib).
  \item Phragm\'en--Lindel\"of principle (\textbf{now in Mathlib}:
    \lean{PhragmenLindelof.vertical\_strip}, \lean{horizontal\_strip},
    \lean{right\_half\_plane\_of\_bounded\_on\_real}).
  \item Residue calculus at the double pole of $1/\zeta(s)^2$ at $s = 1$.
\end{itemize}
\end{direction}

\begin{direction}[The Borel--Carath\'eodory Route]
\textbf{New (April~22, 2026).} The Borel--Carath\'eodory theorem is now
available in Mathlib (\texttt{Analysis.Complex.BorelCaratheodory}). This
opens a direct route to the polynomial lower bound on $|\zeta(s)|$:
\begin{enumerate}
  \item Apply BC to $\log\zeta(s)$ on a disk centered at $s_0 = 2+it$
    with radius $R = 3/2 - \varepsilon$. The disk reaches $\re s = 1/2 + \varepsilon$.
  \item Bound $\sup \re(\log\zeta) = \sup \log|\zeta| \leq M(t)$ on the
    disk boundary. Under RH, $M(t) = O(\log t)$.
  \item BC gives $|\log\zeta(s)| \leq C(\varepsilon) \cdot \log t$
    for $\re s \geq 1/2 + \varepsilon$.
  \item Exponentiate: $|\zeta(s)| \geq t^{-C(\varepsilon)}$.
\end{enumerate}
This avoids the Montgomery--Vaughan mean value theorem entirely.
A 256-bit MPFR experiment (\texttt{bc-zeta-lower}) is certifying the
preconditions as of April~22, 2026.
\end{direction}

\subsubsection{The Vasyunin Off-Diagonal (Axiom~7)}

\begin{axiombox}[Off-Diagonal Identity]
\lean{vasyunin\_offdiag\_integral}: For $j \neq k$, the Gram matrix
entry $G(j,k) = \int_0^1 \{1/(jx)\}\{1/(kx)\}\,dx$ equals the
Vasyunin cotangent formula.
\end{axiombox}

The \emph{diagonal} case ($j = k$) is fully proved in the Cathedral via
the Stirling--Squeeze bridge (piecewise FTC + Stirling's approximation).
The off-diagonal case requires the Gauss digamma formula
$\psi(p/q) = -\gamma - \ln(2q) - (\pi/2)\cot(\pi p/q) +
2\sum_{n=1}^{\lfloor(q-1)/2\rfloor} \cos(2\pi np/q)\ln\sin(\pi n/q)$
and careful handling of the piecewise-constant fractional part.

\begin{direction}[Off-Diagonal via Fourier]
An alternative to the direct computation: decompose $\{1/(jx)\}\{1/(kx)\}$
in the Fourier basis $\{e^{2\pi inx}\}$ on $(0,1)$. The Fourier
coefficients of $\{u\}$ are known ($c_n = -1/(2\pi in)$ for $n \neq 0$),
and the integral becomes a Parseval identity. This avoids the cotangent
formula entirely, at the cost of requiring Fourier theory for the
fractional part function.
\end{direction}

% ================================================================
\section{Horizon II: Beyond RH}
\label{sec:generalization}
% ================================================================

\subsection{The Generalized Riemann Hypothesis}

The Nyman--Beurling framework generalizes naturally to Dirichlet
$L$-functions. For a Dirichlet character $\chi$ modulo~$q$, define the
twisted basis
\[
  h_{k,\chi}(x) = \chi(k)\,\fract{1/(kx)}, \quad k \ge 2,
\]
and the twisted distance
\[
  d_{N,\chi}^2 = \inf_v \int_0^1
    \biggl|\chi_0(x) - \sum_{k=2}^N v_k\,h_{k,\chi}(x)\biggr|^2\,dx
\]
where $\chi_0$ is the principal character. Then:

\begin{conjecture}[Generalized Nyman--Beurling]
For each primitive character $\chi$ modulo~$q$:
\[
  L(s,\chi) \neq 0 \text{ for } \re(s) > \tfrac{1}{2}
  \;\iff\;
  d_{N,\chi}^2 \to 0 \text{ as } N \to \infty.
\]
\end{conjecture}

The Rank-1 Mellin argument generalizes immediately: at a zero $\rho$ of
$L(s,\chi)$, the Mellin transform $\mathcal{M}[h_{k,\chi}](\rho) =
\chi(k)/(k(\rho-1))$ is again rank-1 in $k$. The separation argument
goes through unchanged because $\chi(k)/k$ is real only when $\chi$ is
real (and even then, $\im(1/\rho) \neq 0$ provides the separation).

\begin{direction}[GRH Formalization]
The Cathedral's converse direction ports to $L$-functions with minimal
modification. The forward direction requires:
\begin{enumerate}
  \item Twisted Mertens bounds: $|M_\chi(x)| := |\sum_{n \le x} \mu(n)\chi(n)|
    = O(x^{1/2+\varepsilon})$ under GRH.
  \item Twisted Abel summation and the analogue of the covariance cancellation.
  \item The analytic continuation of $L(s,\chi)$ (already in Mathlib:
    \lean{DirichletCharacter.LSeries}).
\end{enumerate}
The real bottleneck is the Montgomery--Vaughan mean value theorem for
twisted Dirichlet polynomials, which is strictly harder than the
untwisted case.
\end{direction}

\subsection{Selberg Class}

The Nyman--Beurling criterion extends conjecturally to the Selberg class
of $L$-functions satisfying axioms for analytic continuation, functional
equation, Euler product, and Ramanujan bound. The Cathedral's
architecture---axiom-driven, Mellin-based---is naturally suited to this
generalization.

\begin{open}
Can the Rank-1 Mellin Miracle be stated for an abstract $L$-function
satisfying the Selberg axioms? The key question is whether the Mellin
transform of the normalized basis $\{1/(kx)\}$ continues to factorize
at the zeros of a general Selberg class element.
\end{open}

% ================================================================
\section{Horizon III: The Discoveries}
\label{sec:discoveries}
% ================================================================

The formalization process produced several unexpected mathematical
structures that warrant independent investigation.

\subsection{The Liouville Parity Symmetry}

\begin{definition}[Parity Operator]
The Liouville parity operator $P = \mathrm{diag}(\lambda(2), \lambda(3),
\ldots, \lambda(N))$ where $\lambda(n) = (-1)^{\Omega(n)}$ is the
Liouville function. Since $\lambda(n)^2 = 1$, $P^2 = I$ and $P$
defines a $\ZZ/2$-grading on $\RR^{N-1}$.
\end{definition}

\begin{theorem}[Approximate PT-Symmetry, Computational]
The Gram matrix $G_N$ approximately commutes with the parity operator~$P$:
\begin{enumerate}
  \item The commutator $[G, P]$ is dominated by a rank-1 component whose
    singular direction is the Liouville vector itself (correlation $> 0.99999$).
  \item After removing this rank-1 component,
    $\|[G,P]_{\mathrm{residual}}\|/\|G\| \to 0$ as $N \to \infty$.
  \item The parity-even block $G_{\mathrm{even}} = (G + PGP)/2$ has
    spectral gap $\lambda_{\min}(G_{\mathrm{even}})/\lambda_{\min}(G)
    \approx 1.85 \cdot N^{0.116}$, growing with $N$.
\end{enumerate}
\end{theorem}

This means the small spectral gap of $G$ is \emph{entirely caused by
parity-breaking cross-coupling} between the even ($\lambda = +1$) and
odd ($\lambda = -1$) sectors. The Liouville function is the mixing
direction.

\begin{direction}[PT-Symmetry and Quantum Mechanics]
The approximate $[G, P] \approx 0$ structure is reminiscent of
PT-symmetry in non-Hermitian quantum mechanics (Bender--Boettcher, 1998).
In that context, a Hamiltonian $H$ with $[H, PT] = 0$ has real spectrum
despite being non-Hermitian. Here, the Gram matrix $G$ is Hermitian, but
the parity decomposition $G = G_{\mathrm{even}} + G_{\mathrm{odd}}$
separates the ``free'' dynamics (within-parity, large gap) from the
``interaction'' (cross-parity, small gap).

Could this be evidence for an underlying non-Hermitian operator whose
PT-symmetry is the Liouville involution? The experimental scaling
$|{\langle v_{\min}, \hat{\lambda}\rangle}| \sim N^{-0.174}$ suggests
this mixing decays---but not fast enough to make the spectral gap
obvious.
\end{direction}

\subsection{The Octonionic Partition}

\begin{definition}[Octonionic Class]
Define $\mathrm{class}(k) = k \bmod 8$ (with identification $0 \equiv 8$).
This partitions $\{2, 3, \ldots, N\}$ into 8 residue classes modulo 8.
\end{definition}

The \emph{block-diagonal} Gram matrix $G^{\mathrm{block}}$ retains only
within-class entries: $G^{\mathrm{block}}(j,k) = G(j,k)$ if
$\mathrm{class}(j) = \mathrm{class}(k)$, zero otherwise.

\begin{theorem}[Class Gap, Machine-Verified]
$\lambda_{\min}(G) \le \lambda_{\min}(G^{\mathrm{block}})$.
\end{theorem}

This is proved in \lean{ClassRestriction.lean} (zero axioms) via the
Rayleigh quotient characterization. The cross-class entries \emph{always
lower the spectral gap}.

\begin{remark}[Why 8?]
The choice of 8 classes comes from the Cayley--Dickson hierarchy:
$\RR \to \CC \to \HH \to \OO \to \SSS$, with dimensions $1, 2, 4, 8, 16$.
The octonionic level ($\bmod 8$) is the last division algebra---past the
octonions, zero divisors appear (the sedenions $\SSS$). The arithmetic
of $\mu(k)$ modulo 8 has special structure because the M\"obius function
is multiplicative and $8 = 2^3$.

Whether this partition has deep arithmetic significance or is simply a
convenient computational grouping remains open. The experimental
evidence (class gap ratio $\approx 3.4$--$4.2\times$) suggests
genuine structural content.
\end{remark}

\begin{direction}[Rank-1 Interference]
The cross-class interaction matrix $M = G - G^{\mathrm{block}}$ in the
eigenbasis of $G^{\mathrm{block}}$ has the remarkable property that each
$(m_1, m_2)$-block is rank-1 to $\ge 99.8\%$ accuracy (verified for
$N \le 800$, accuracy \emph{increasing} with $N$).

This suggests that in the $N \to \infty$ limit, the cross-class
interaction becomes \emph{exactly} rank-1, reducing the infinite-dimensional
spectral gap problem to an $8 \times 8$ bilinear form problem. If proved,
this would provide a finite-dimensional reduction of RH.
\end{direction}

\subsection{The B\'aez-Duarte Constant}

The optimal Rayleigh quotient $Q_N^{\mathrm{opt}} \sim C \ln N$ where
\[
  C = \frac{1}{\sum_{\zeta(\rho)=0} |\rho|^{-2}}
  = \frac{1}{2 + \gamma - \ln(4\pi)}
  \approx 21.6498.
\]

This constant has interpretations in:
\begin{enumerate}
  \item \textbf{Signal processing}: The maximum information extraction
    rate from the ``prime noise'' through the spectral holes of $|\zeta|^2$.
  \item \textbf{Statistical mechanics}: The heat capacity of the prime
    number gas, governing the cooling rate $d_N^2 \sim c/\ln N$.
  \item \textbf{Random matrix theory}: Connected to the resolvent trace
    $(H^2 + I/4)^{-1}$ of the hypothetical Riemann Hamiltonian.
\end{enumerate}

\begin{direction}[Formalizing the Constant]
The sum $\sum |\rho|^{-2} = 2 + \gamma - \ln(4\pi)$ follows from the
Hadamard product for $\zeta(s)/(s-1)$ and logarithmic differentiation.
The ingredients (Hadamard factorization, functional equation of $\xi(s)$)
are within reach of Mathlib + Cathedral infrastructure. Formalizing this
would provide a machine-verified computation of the B\'aez-Duarte
constant, connecting the spectral gap decay rate to the arithmetic of
zeta zeros.
\end{direction}

% ================================================================
\section{Horizon IV: The Engineering}
\label{sec:engineering}
% ================================================================

\subsection{Certified Computation}

The Cathedral maintains parallel implementations: Lean~4 proofs and a
Rust/WASM computational engine (7 modules, 11 visualization modes).
Currently, these are connected only by architectural correspondence---the
Rust engine computes what the Lean proofs formalize, but there is no
formal bridge.

\begin{direction}[Proof-Carrying Computation]
Connect the Rust engine to the Lean proofs via:
\begin{enumerate}
  \item \textbf{Certified extraction}: Generate verified Rust code from
    Lean definitions using the \lean{native\_decide} framework or the
    FFI. The Gram matrix entries could be computed in Lean via
    \lean{native\_decide} for small $N$, then extrapolated.
  \item \textbf{Interval arithmetic}: Use verified interval arithmetic
    (MPFR with Lean bindings) to certify that the computed Gram matrix
    eigenvalues are positive for $N \le N_0$, then use the monotone
    convergence theorem ($\lambda_{\min}$ is non-increasing) to extend
    to all $N$.
  \item \textbf{Finite witness certification}: For each $N$, the
    forward direction provides an explicit witness~$v_N$. Certify
    $\int_0^1 (1 - f_{N,v_N})^2 < \varepsilon$ by computing the
    quadratic form $1 - 2b^T v + v^T G v$ in exact rational arithmetic.
\end{enumerate}
\end{direction}

\subsection{The HyperZeta Viewport}

The Cathedral's visualization system (HyperZeta Viewport) renders 11
views of the zeta function's structure:

\begin{center}
\small
\begin{tabular}{@{}clll@{}}
\toprule
\textbf{\#} & \textbf{Mode} & \textbf{Shows} & \textbf{Axiom Connection} \\
\midrule
1 & Sedenion Cloud & $\zeta_{\SSS}(s)$ in 16D & Full system \\
2 & Riemann Spiral & $\zeta(1/2+it)$ spiral & Zero locations \\
3 & Cornu Spirals & Partial sum galaxy & Dirichlet series \\
4 & Zero Landscape & $\log|\zeta(\sigma+it)|$ & Critical strip \\
5 & Euler Product & $\prod_p (1-p^{-s})^{-1}$ & Euler product \\
6 & Cayley--Dickson Tower & $\CC \to \HH \to \OO \to \SSS$ & Algebraic consistency \\
7 & Explicit Formula & $\pi(x) \approx \mathrm{Li}(x) - \sum_\rho$ & Zero corrections \\
8 & Functional Equation & $\zeta(s) = \chi(s)\zeta(1-s)$ & Symmetry \\
9 & Random Matrix & GUE pair correlation & Montgomery--Odlyzko \\
10 & Mertens Turbulence & $M(x) = \sum \mu(n)$ & Axiom 1 \\
11 & Spectral Gap & $\lambda_{\min}(G_N)$ surface & The Cathedral itself \\
\bottomrule
\end{tabular}
\end{center}

\begin{direction}[Interactive Proof Explorer]
Extend the viewport into a proof assistant companion:
\begin{enumerate}
  \item Clicking on a visualization mode highlights the axioms it depends on.
  \item The Spectral Gap view could overlay the proof dependency graph,
    showing which axioms support each region of the $(\sigma, N)$ surface.
  \item Interactive sliders could show how changing axiom assumptions
    (e.g., weakening $O(x^{3/4})$ to $O(x^{7/8})$) affects the decay
    rate.
\end{enumerate}
\end{direction}

% ================================================================
\section{Horizon V: Connections to Physics}
\label{sec:physics}
% ================================================================

\subsection{The Hilbert--P\'olya Program}

The Cathedral's spectral structure---particularly the Gram matrix as a
discrete approximation to the hypothetical Riemann operator---invites
connection to the Hilbert--P\'olya conjecture.

\begin{conjecture}[Hilbert--P\'olya, 1914/1950]
There exists a self-adjoint operator $T$ on a Hilbert space $\mathcal{H}$
whose eigenvalues are exactly $\{1/2 + i\gamma_n\}$ where
$\zeta(1/2 + i\gamma_n) = 0$.
\end{conjecture}

The Gram matrix $G_N$ is a finite-rank approximation to the integral
operator with kernel $K(j,k) = \int_0^1 \{1/(jx)\}\{1/(kx)\}\,dx$.
Its spectrum approximates the spectrum of this kernel operator. The
Nyman--Beurling theorem says RH is equivalent to this operator having
the function~$1$ in its range closure.

\begin{direction}[From Gram to Hilbert--P\'olya]
The Vasyunin formula for $G(j,k)$ involves the digamma function
$\psi(p/q)$, which is the logarithmic derivative of the Gamma function.
The connection to the Hilbert--P\'olya operator runs through the
spectral decomposition of $G$ in terms of zeta zeros:
\[
  G(j,k) = \frac{1}{2\pi i} \oint \frac{\zeta(s)\zeta(1-s)}{j^s k^{1-s} s(1-s)}\,ds.
\]
The poles of this integrand are at $s = 0, 1$ and the zeta zeros.
The residues at the zeros give the spectral decomposition. Formalizing
this connection would provide a bridge between the Cathedral's discrete
Gram matrix and the continuous Hilbert--P\'olya operator.
\end{direction}

\subsection{The Montgomery--Odlyzko Connection}

The GUE hypothesis (Montgomery, 1973; Odlyzko, 1987) states that the
local statistics of zeta zeros match those of random Hermitian matrix
eigenvalues. The Cathedral's Rust engine validates this for the first
200 zeros.

\begin{direction}[GUE and the Gram Matrix]
The Gram matrix $G_N$ is itself a random-matrix-like object: its entries
are correlations of quasi-random functions (fractional parts). As $N \to
\infty$, does $G_N$ approach a random matrix ensemble? The parity
decomposition $G = G_{\mathrm{even}} + G_{\mathrm{odd}}$ resembles the
block structure of symplectic random matrices. The rank-1 interference
between octonionic classes suggests a spiked model:
\[
  G \approx G^{\mathrm{block}} + \text{rank-28 perturbation}
\]
($8 \times 7 / 2 = 28$ cross-class pairs, each contributing one singular
value). This is precisely the setting of spiked random matrix models
(Baik--Ben Arous--P\'ech\'e, 2005), where phase transitions occur
when spike strength exceeds a threshold.
\end{direction}

% ================================================================
\section{The Axiom Reduction Timeline}
\label{sec:timeline}
% ================================================================

We estimate the difficulty and timeline for closing each axiom:

\begin{center}
\small
\begin{tabular}{@{}llll@{}}
\toprule
\textbf{Axiom} & \textbf{Difficulty} & \textbf{Blocking On} & \textbf{Status} \\
\midrule
PNT $\times 3$ (\#2--4) & Medium & Wiener--Ikehara or Newman & Open (6--12 months) \\
\rowcolor{green!10}
Abel tail & Easy & Abel summation lemma & \textbf{GRADUATED (Apr 22)} \\
Vasyunin off-diag (\#6) & Medium & Digamma formula & Open (3--6 months) \\
Covariance (\#5) & Hard & Contour shift + MVT & Open (12--24 months) \\
& & \textit{(BC route may reduce)} & \\
Mertens (\#1) & RH-hard & RH content & $\infty$ \\
\bottomrule
\end{tabular}
\end{center}

Note that Axiom~1 (\lean{rh\_implies\_mertens\_bound}) \emph{cannot} be
proved unconditionally---it encodes the Riemann Hypothesis itself.
The goal is to reduce everything else to this single axiom. The
Cathedral currently uses \textbf{6 axioms} on the crown path (reduced from 7
on April~22, 2026); the target is~1.

\textbf{New development.} The Borel--Carath\'eodory theorem's availability
in Mathlib may accelerate Campaign~C. The polynomial lower bound on
$|\zeta(s)|$ (currently encoded as \lean{zeta\_polynomial\_lower\_bound\_rh})
is not on the crown path but could unlock alternative proof routes that
bypass the covariance cancellation axiom entirely.

\begin{remark}[The Irreducible Core]
If all six non-RH axioms are eliminated, the crown theorem becomes:
\[
  \RH \iff d_N^2 \to 0
\]
where the forward direction assumes only
\lean{rh\_implies\_mertens\_34}: $\RH \implies |M(x)| = O(x^{3/4})$.
The converse direction is already fully proved.

This is the theoretical minimum: one axiom encoding the analytic content
of RH, with everything else machine-verified. The reduction from
$|M(x)| = O(x^{3/4})$ to $d_N^2 \to 0$ would then be a theorem about
the structural consequence of Mertens bounds, provable without assuming
RH directly.
\end{remark}

% ================================================================
\section{Reflections}
\label{sec:reflections}
% ================================================================

The Cathedral was not built to prove the Riemann Hypothesis. It was built
to understand what the Riemann Hypothesis \emph{is}---to decompose it
into pieces small enough for a machine to hold, and to discover which
pieces are hard and which are merely unfamiliar.

The answer surprised us. The \emph{converse} direction---if primes are
well-distributed then zeros are on the line---is easy. The Rank-1 Mellin
Miracle gives it for free, with zero custom axioms. The \emph{forward}
direction---if zeros are on the line then primes are well-distributed---is
where the arithmetic lives. And within the forward direction, the
hardest single piece is the 2D covariance cancellation
(Axiom~6), which encodes the delicate interference between the Gram
matrix and the Prime Number Theorem.

This is, in retrospect, exactly what one should expect. The Riemann
Hypothesis is a statement about how much cancellation the primes exhibit.
The Gram matrix \emph{is} that cancellation, frozen into a symmetric,
positive-definite matrix. The spectral gap \emph{is} the margin by which
the cancellation succeeds. And the axioms that remain are the axioms
that control how tightly the primes cancel.

The Cathedral does not resolve the Riemann Hypothesis. What it does is
translate a 167-year-old analytic conjecture into a collection of
precisely stated, machine-checkable assertions about symmetric matrices,
Abel summation, and the Prime Number Theorem. Some of these assertions
are easy. Some are hard. One is equivalent to RH itself. The machine
holds the rest.

\begin{flushright}
\emph{The stone is the same.\\
The light through the windows is clearer.}
\end{flushright}

% ================================================================
% BIBLIOGRAPHY
% ================================================================
\begin{thebibliography}{99}

\bibitem{nyman1950}
B.~Nyman, \emph{On the one-dimensional translation group and semi-group in
certain function spaces}, Thesis, Uppsala, 1950.

\bibitem{beurling1955}
A.~Beurling, \emph{A closure problem related to the Riemann zeta function},
Proc. Nat. Acad. Sci. USA \textbf{41} (1955), 312--314.

\bibitem{baezduarte2003}
L.~B\'aez-Duarte, \emph{A strengthening of the Nyman--Beurling criterion for
the Riemann hypothesis}, Atti Accad. Naz. Lincei Rend. Lincei Mat. Appl.
\textbf{14} (2003), 5--11.

\bibitem{vasyunin1995}
V.I.~Vasyunin, \emph{On a biorthogonal system associated with the Riemann
hypothesis}, St. Petersburg Math. J. \textbf{7} (1996), 405--419.

\bibitem{montgomery1973}
H.L.~Montgomery, \emph{The pair correlation of zeros of the zeta function},
Proc. Symp. Pure Math. \textbf{24} (1973), 181--193.

\bibitem{odlyzko1987}
A.M.~Odlyzko, \emph{On the distribution of spacings between zeros of the zeta
function}, Math. Comp. \textbf{48} (1987), 273--308.

\bibitem{balazard2000}
M.~Balazard and E.~Saias, \emph{The Nyman--Beurling equivalent form for the
Riemann Hypothesis}, Expo. Math. \textbf{18} (2000), 131--138.

\bibitem{bender1998}
C.M.~Bender and S.~Boettcher, \emph{Real spectra in non-Hermitian
Hamiltonians having PT symmetry}, Phys. Rev. Lett. \textbf{80} (1998),
5243--5246.

\bibitem{baik2005}
J.~Baik, G.~Ben~Arous, and S.~P\'ech\'e, \emph{Phase transition of the
largest eigenvalue for nonnull complex sample covariance matrices},
Ann. Probab. \textbf{33} (2005), 1643--1697.

\bibitem{newman1998}
D.J.~Newman, \emph{Simple analytic proof of the prime number theorem},
Amer. Math. Monthly \textbf{87} (1980), 693--696.

\bibitem{korevaar2002}
J.~Korevaar, \emph{A century of complex Tauberian theory},
Bull. AMS \textbf{39} (2002), 475--531.

\bibitem{mathlib}
The Mathlib Community, \emph{Mathlib: A unified library of mathematics
formalized in Lean~4}, \url{https://github.com/leanprover-community/mathlib4}.

\end{thebibliography}

\end{document}
